\section{Determination and sufficient source-occurrence theorems}\label{sec:occurrence}
\subsection{Off-diagonal data leave exactly three local constants}
Suppose B is a continuous Hermitian form on $\cD^0$,
stationary under $U_t$ and invariant under reflection.
Assume its mixed pairing on $\T h,\T k$ agrees with Q
whenever h and k have disjoint supports. This hypothesis
includes the prime shifts and the archimedean off-diagonal
kernel, with their differentiated pullbacks.

\begin{theorem}[Three-anchor determination]\label{thm:determination}
If also $B[f_\lambda]=O_h(\log|\lambda|)$ for the normalized
modulations \eqref{eq:modulation}, then
\begin{equation}\label{eq:local-ambiguity}
B(\T h,\T k)-Q(\T h,\T k)
=a_0\ip h k_2+a_1\ip{h'}{k'}_2+a_2\ip{h''}{k''}_2
\end{equation}
for three real constants. Equality on three anchors
$\T h_r$, where $h_r(x)=r^{-1/2}h(x/r)$, any nonzero
$h\in\cD$, and $r=1,2,3$, forces $B=Q$ everywhere.
\end{theorem}
\begin{proof}
Pull the difference back to $\cD$ through the topological
isomorphism $\T$. The distribution kernel theorem and
stationarity give a convolution kernel. Disjoint-support
vanishing makes that kernel supported at zero, hence it is
a finite linear combination of derivatives of $\delta$.
Hermiticity makes its Fourier polynomial real; reflection
makes it even. For
$f_\lambda=\lambda^{-2}\T(e^{\ii\lambda x}h)$, a term of
degree d has leading size $\lambda^{d-4}\norm h^2$.
The logarithmic bound and Proposition~\ref{prop:highfreq}
therefore exclude $d>4$. This proves
\eqref{eq:local-ambiguity}.
At the three scaled anchors the coefficients multiply
$\norm h^2,r^{-2}\norm{h'}^2,r^{-4}\norm{h''}^2$.
The determinant is
$-5\norm h^2\norm{h'}^2\norm{h''}^2/54$, which is nonzero.
\end{proof}
Three constants are necessary in this argument. Without
reflection the Hermitian polynomial may have every degree
from zero through four, leaving five real constants.
Symmetry, growth, and finitely many anchors \emph{without}
the off-diagonal identity are inadequate: a real even Schwartz
multiplier can be chosen to vanish on all those finitely
many anchor functionals yet change the form.

One useful necessary test is visible after pullback. At
each nonzero prime shift the kernel of $Q(\T h,\T k)$
contains
\begin{equation}
-\frac{\Lambda(a)}{\sqrt a}
\T_r^2\{\delta(r-\log a)+\delta(r+\log a)\},
\qquad \T_r^2=\partial_r^4-\tfrac12\partial_r^2+\tfrac1{16}.
\end{equation}
A candidate with a smooth off-diagonal mixed kernel cannot
supply these singularities. Negative off-diagonal coefficients
alone are not an obstruction to positivity: squared shift
differences supply such terms together with positive contact.

\subsection{Compactness turns verified finite feasibility into occurrence}
\begin{theorem}[Compact-source occurrence criterion]\label{thm:compact-occurrence}
Let $\cK\Subset\cH_{\rm phys}$ be an injectively and compactly
embedded Hilbert source space. Choose a countable complex-rational
linear test core in $\cD^0$, dense on every enclosing compact
support interval. Prescribe for each coordinate f a bound
$\norm{Jf}_{\cK}\le C_Ip_I(f)$, using fixed continuous test
seminorms on each interval I. Suppose every finite collection
of the desired source relations is feasible by actual vectors
in these same balls. The relations must be closed in the
strong physical topology and include complex-rational
linearity and the required mixed pairings on the test core.
Then there is a continuous global source law satisfying all
these relations and the same bounds.
\end{theorem}
\begin{proof}
Each $\cK$ ball has compact image in $\cH_{\rm phys}$.
Its image is closed: a strong physical limit of bounded
$\cK$ vectors has a weak $\cK$ subsequence, whose image
equals that strong limit. Take the product of these compact
coordinate balls. Every imposed closed relation cuts out a
closed subset. Finite feasibility is exactly the finite
intersection property, so a common assignment exists.
For each support interval the seminorm bound extends the
linear assignment from the dense core to the full test
space. Rational complex linearity and continuity give
complex linearity. Compatible approximations in larger
intervals give one global law; physical strong continuity
passes all mixed pairings to the limit. Weak compactness
in $\cK$ preserves the stated source bounds.
\end{proof}
Unbounded operator relations require their actual closed
graph and sufficient coordinate bounds, rather than a formal
symbolic identity. A high-frequency estimate needed for
Theorem~\ref{thm:determination} must likewise be imposed
uniformly by a common continuous seminorm; asymptotic
agreement on a countable list does not imply it.
A useful special case is a sequence of full source maps
with uniform compact-source bounds and convergent mixed
data on a dense core; a diagonal compactness argument
then gives the same conclusion.

The finite boundary model supplies the spaces $\cK_r$ in
\eqref{eq:compact-tail}, but the raw character packets fail
their bounds. More fundamentally, feasibility of the
arithmetic relations with these bounds has not been
established. Positivity of separately assembled finite
matrices does not establish the theorem's physical
finite-feasibility hypothesis.

\subsection{Optional continuum OS and phase-space hypotheses}
Assume a reconstructed vacuum theory with local algebras
$\mathcal A(O)$, vacuum $\Omega_{\rm vac}$, and physical
Hamiltonian $H_{\rm phys}\ge0$. A mass gap means
$H_{\rm phys}\ge\Delta(1-P_{\rm vac})$. It controls infrared
excitations, but does not imply compactness of arbitrary
ultraviolet source families.

A natural additional phase-space hypothesis is that, for
bounded O and $\beta>0$,
\begin{equation}
\Theta_{\beta,O}:\mathcal A(O)\longrightarrow\cH_{\rm OS},
\qquad A\longmapsto e^{-\beta H_{\rm phys}}A\Omega_{\rm vac}
\end{equation}
maps the operator-norm unit ball into a relatively compact
set. Nuclearity is a familiar stronger condition
\cite{buchholz-wichmann}; it is an explicit extra assumption
here, not a consequence claimed from the mass gap.

For each finite M the set
\begin{equation}
\mathcal K(O,M)=
\{e^{-\beta H_{\rm phys}}A\Omega_{\rm vac}:
                A\in\mathcal A(O),\ \norm A\le M\}
\end{equation}
is actually compact. Relative compactness is assumed.
For closedness, take a norm-convergent sequence of images,
an ultraweak convergent subnet of their bounded operators,
and test against arbitrary vectors. The vector functionals
are normal, so the norm limit is the image of the ultraweak
limit operator in $\mathcal A(O)$ with the same bound.
There is also the elementary energy-tail estimate
\begin{equation}
\norm{\mathbf1_{[E,\infty)}(H_{\rm phys})
          e^{-\beta H_{\rm phys}}A\Omega_{\rm vac}}
\le Me^{-\beta E}.
\end{equation}

Consequently the preceding finite-intersection proof applies
with these compact prepared-vector sets. If every finite
set of required arithmetic/source relations is realizable
with fixed local regions $O_I$ and fixed operator bounds
$M_f\le C_Ip_I(f)$, it produces a global actual source
$Jf=e^{-\beta H_{\rm phys}}A_f\Omega_{\rm vac}$ with those
relations. If the vacuum is separating for the relevant
local algebra, injectivity makes $A_f$ unique and transfers
vector linearity to operator linearity; otherwise only the
vector law is needed. Reeh--Schlieder density by itself
does not provide the uniform bounds or arithmetic
feasibility. No such feasibility is proved here, even
when all OS axioms, a mass gap, and phase-space compactness
are granted.
